% Greek abstract (required by the department)
\begin{abstract}
Η πρόβλεψη ακμών σε γράφους με κειμενικά χαρακτηριστικά κόμβων αποτελεί το πρόβλημα
εξαγωγής συμπεράσματος για την ύπαρξη ακμής μεταξύ δύο κόμβων, λαμβάνοντας υπόψη
τόσο την τοπολογία του γράφου όσο και το κειμενικό περιεχόμενο των κόμβων. Ενώ
πρόσφατες εργασίες στο σύνολο αξιολόγησης \en{DSAA 2023} της \en{Wikipedia} κατέδειξαν
ότι καθαρά σημασιολογικές προσεγγίσεις βασισμένες σε προ-εκπαιδευμένα γλωσσικά
μοντέλα μπορούν να επιτύχουν σχεδόν τέλεια ακρίβεια, το τετραγωνικό τους κόστος
συμπερασμού τα καθιστά μη πρακτικά σε κλίμακα, ενώ η γενίκευσή τους σε ακμές εκτός
κατανομής παραμένει αδιερεύνητη.

Η παρούσα διπλωματική εργασία παρουσιάζει το \textbf{CascadeLP}, ένα πολυεπίπεδο
πλαίσιο πρόβλεψης ακμών που δρομολογεί ζεύγη κόμβων μέσω προοδευτικά πιο δαπανηρών
μοντέλων βάσει ορίων αξιοπιστίας. Το Επίπεδο~1 εφαρμόζει δομικές ευρεσιτεχνίες
(\en{Common Neighbours}, \en{Adamic-Adar}, \en{Node2Vec}) για ταχεία, τοπολογικά
ενήμερη βαθμολόγηση. Το Επίπεδο~2 εφαρμόζει έναν ελαφρύ σημασιολογικό ταξινομητή
(χαρακτηριστικά \en{TF-IDF} + \en{SVM} ή συχνότητες μερών του λόγου + Τυχαίο Δάσος)
για τον χειρισμό κόμβων χωρίς δομικούς γείτονες. Το Επίπεδο~3 εφαρμόζει ένα μοντέλο
\en{Sentence-RoBERTa} αποκλειστικά για ζεύγη στα οποία τα προηγούμενα επίπεδα
εμφανίζουν αβεβαιότητα. Αξιολογούμε το \en{CascadeLP} στο σύνολο δεδομένων
\en{DSAA 2023} τόσο με τυπική όσο και με χρονική διαίρεση, και αναφέρουμε ακρίβεια,
ρυθμό επεξεργασίας και ποσοστά κλήσης ανά επίπεδο. Τα αποτελέσματά μας δείχνουν ότι
το \en{CascadeLP} εξισώνει την ακρίβεια του πλήρους μετασχηματιστή ενεργοποιώντας τον
αποκλειστικά σε ένα μικρό κλάσμα ζευγών, παρέχοντας μια συγκεκριμένη ανάλυση
συμβιβασμού αποτελεσματικότητας-ακρίβειας που προηγούμενες εργασίες δεν έχουν
επιχειρήσει.

    \begin{keywords}
    Πρόβλεψη ακμών, γράφοι με κειμενικά χαρακτηριστικά, νευρωνικά δίκτυα γράφων,
    γλωσσικά μοντέλα, βαθμιδωτή συμπερασματολογία, εκκίνηση εκ του μηδενός,
    \en{Wikipedia}, \en{DSAA 2023}
    \end{keywords}
\end{abstract}


% English abstract
\begin{abstracteng}
\en{
Link prediction in text-attributed graphs (TAGs) is the problem of inferring whether an
edge exists between two nodes given both the graph topology and the textual content of
the nodes. While recent work on the DSAA 2023 Wikipedia benchmark demonstrated that
pure semantic approaches based on pre-trained language models can achieve near-perfect
accuracy, their quadratic inference cost makes them impractical at scale, and their
generalisation to out-of-distribution edges remains unstudied.

This thesis presents \textbf{CascadeLP}, a tiered link prediction framework that routes
node pairs through progressively expensive models based on confidence thresholds.
Tier~1 applies structural heuristics (Common Neighbours, Adamic-Adar, Node2Vec) for
fast, topology-aware scoring. Tier~2 applies a lightweight semantic classifier
(TF-IDF + SVM or POS-tag features + Random Forest) to handle cold-start nodes with no
structural neighbours. Tier~3 applies a fine-tuned Sentence-RoBERTa model only to pairs
where the earlier tiers are uncertain. We evaluate CascadeLP on the DSAA 2023 dataset
under both standard and temporal splits, and report accuracy, throughput, and per-tier
call rates. Our results show that CascadeLP matches full-transformer accuracy while
invoking the transformer on only a small fraction of pairs, providing a concrete
efficiency-accuracy trade-off analysis that prior work has not attempted.
}

    \begin{keywordseng}
    \en{Link prediction, text-attributed graphs, graph neural networks, language models,
    cascaded inference, cold-start, Wikipedia, DSAA 2023}
    \end{keywordseng}
\end{abstracteng}
